\vspace{-0.3em}
\textit{ShiftRows} es \begin{Revision}la primera transformación que se efectúa dentro de la capa de difusión\end{Revision}, cuyo objetivo es propagar redundancias en el texto claro a través del texto cifrado, de manera que no se puedan apreciar dichas redundancias en el texto cifrado. El algoritmo de \textit{ShiftRows} consiste en tomar cada fila $i$ del \textit{estado} y desplazar sus bytes $i$ unidades a la izquierda, tal y como se muestra en la Figura \ref{fig:ShiftRows}. A dicho desplazamiento se le denomina rotación cíclica hacia la izquierda/derecha en $i$ unidades, pues los elementos situados en los extremos se desplazan al extremo opuesto (de forma cíclica).
\vspace{-0.2em}
\begin{figure}[H] 
  \centering 
  
  \begin{tikzpicture}[baseline=(current bounding box.center), >=Stealth]
    
    % 1. Definición de colores pastel sutiles para cada fila
    \definecolor{filaAzul}{RGB}{235, 243, 250}
    \definecolor{filaVerde}{RGB}{236, 248, 236}
    \definecolor{filaAmarilla}{RGB}{254, 249, 231}
    \definecolor{filaRoja}{RGB}{253, 235, 235}

    % 2. Matriz de la izquierda (Estado Inicial)
    \matrix (A) [
      matrix of math nodes,
      nodes={draw, minimum size=0.9cm, anchor=center},
      column sep=-\pgflinewidth,
      row sep=-\pgflinewidth
    ] at (0,0) {
      |[fill=filaAzul]| a_0     & |[fill=filaAzul]| a_4     & |[fill=filaAzul]| a_8     & |[fill=filaAzul]| a_{12} \\
      |[fill=filaVerde]| a_1    & |[fill=filaVerde]| a_5    & |[fill=filaVerde]| a_9    & |[fill=filaVerde]| a_{13} \\
      |[fill=filaAmarilla]| a_2 & |[fill=filaAmarilla]| a_6 & |[fill=filaAmarilla]| a_{10} & |[fill=filaAmarilla]| a_{14} \\
      |[fill=filaRoja]| a_3     & |[fill=filaRoja]| a_7     & |[fill=filaRoja]| a_{11}  & |[fill=filaRoja]| a_{15} \\
    };

    % 3. Matriz de la derecha (Estado Final - ShiftRows)
    \matrix (B) [
      matrix of math nodes,
      nodes={draw, minimum size=0.9cm, anchor=center},
      column sep=-\pgflinewidth,
      row sep=-\pgflinewidth
    ] at (6,0) {
      |[fill=filaAzul]| a_0      & |[fill=filaAzul]| a_4      & |[fill=filaAzul]| a_8     & |[fill=filaAzul]| a_{12} \\
      |[fill=filaVerde]| a_5     & |[fill=filaVerde]| a_9     & |[fill=filaVerde]| a_{13}  & |[fill=filaVerde]| a_1 \\
      |[fill=filaAmarilla]| a_{10} & |[fill=filaAmarilla]| a_{14} & |[fill=filaAmarilla]| a_2  & |[fill=filaAmarilla]| a_6 \\
      |[fill=filaRoja]| a_{15}   & |[fill=filaRoja]| a_3      & |[fill=filaRoja]| a_7     & |[fill=filaRoja]| a_{11} \\
    };

    % 4. Dibujo de las 4 flechas independientes con sus etiquetas de desplazamiento
    \draw[->, thick] (A-1-4.east) -- (B-1-1.west) 
      node[midway, above, font=\scriptsize\sffamily] {Despl. 0 izq.};
      
    \draw[->, thick] (A-2-4.east) -- (B-2-1.west) 
      node[midway, above, font=\scriptsize\sffamily] {Despl. 1 izq.};
      
    \draw[->, thick] (A-3-4.east) -- (B-3-1.west) 
      node[midway, above, font=\scriptsize\sffamily] {Despl. 2 izq.};
      
    \draw[->, thick] (A-4-4.east) -- (B-4-1.west) 
      node[midway, above, font=\scriptsize\sffamily] {Despl. 3 izq.};

  \end{tikzpicture}

  % CAPTION Y ETIQUETA 
  \vspace{-0.2em}
  \caption{Algoritmo \textit{ShiftRows}}
  \label{fig:ShiftRows} % Sirve para citar la imagen en el texto usando \ref{fig:shift_rows_aes}
\end{figure}
\vspace{-0.3em}
Para descifrar se utiliza la operación inversa \textit{InvShiftRows}, que consiste en una rotación cíclica hacia la derecha en $i$ unidades, según se corresponda con la fila $i$.